% =====================================================
% 专题：特殊值与估算专题 (questions.tex)
% =====================================================

% -----------------------------------------------------
% Q1: 反例排除 (函数极值系数)
% -----------------------------------------------------
\newcommand{\qOneStem}{设函数 \(f(x) = x^3 + ax^2 + bx + c\)。若 \(f(x)\) 既有极大值又有极小值，则下列一定成立的是\blank{1.5cm}}
\newcommand{\qOneOptions}{\mcoptions[4]{\(a > 0\)}{\(b < 0\)}{\(a^2 > 3b\)}{\(a + b > 0\)}}
\newcommand{\qOneAnswer}{C}
\newcommand{\qOneSolution}{%
  【解题思路】\\
  本题是判断“一定成立”的结论。使用特殊值（特殊函数模型）进行证伪排它是最快的方法。\\
  第一步：构造符合题干（既有极大值又有极小值）的简单三次函数反例；\\
  第二步：用反例检验各个选项，排除不正确选项；\\
  第三步：进行正面代数验证以锁死答案。\\
  【解析计算】\\
  法一（特殊值反例排除法）：\\
  1. 构造第一个函数 \(f(x) = x^3 - x\)，显然它有极大值和极小值（导数 \(f'(x) = 3x^2 - 1 = 0\) 有两根）。此时 \(a = 0, b = -1\)。代入选项：\\
  对于 A，\(a > 0\) 不成立，排除 A；\\
  对于 D，\(a + b = -1 > 0\) 不成立，排除 D。\\
  2. 构造第二个函数 \(f(x) = x^3 + x^2\)，它也有极值点（导数 \(f'(x) = 3x^2 + 2x = 0\) 有两根）。此时 \(a = 1, b = 0\)。代入选项：\\
  对于 B，\(b < 0\) 不成立，排除 B。\\
  综上，排除 A、B、D，只有 C 正确。\\
  
  法二（正面解析法）：\\
  已知 \(f(x) = x^3 + ax^2 + bx + c\)，求导得：\\
  \(f'(x) = 3x^2 + 2ax + b\)。\\
  因为 \(f(x)\) 既有极大值又有极小值，所以导函数 \(f'(x) = 0\) 必须有两个不同的实数根。\\
  根据一元二次方程根的判别式，得：\\
  \(\Delta = (2a)^2 - 4 \cdot 3 \cdot b > 0 \implies 4a^2 - 12b > 0 \implies a^2 > 3b\)。\\
  故一定成立的是 C。%
}

% -----------------------------------------------------
% Q2: 范围最值 (三角函数最值)
% -----------------------------------------------------
\newcommand{\qTwoStem}{在锐角三角形 \(ABC\) 中，\(\tan A + \tan B + \tan C\) 的最小值为\blank{1.5cm}}
\newcommand{\qTwoOptions}{\mcoptions[4]{\(2\sqrt{3}\)}{\(3\sqrt{3}\)}{\(4\sqrt{3}\)}{\(6\sqrt{3}\)}}
\newcommand{\qTwoAnswer}{B}
\newcommand{\qTwoSolution}{%
  【解题思路】\\
  三角形中最值问题，最对称的情况往往是取得最值时的特殊状态。\\
  第一步：先考虑最对称的等边三角形状态，算出一个合法可取的值，用它来卡住范围边界；\\
  第二步：用该值排除不合理的选项；\\
  第三步：用凸函数性质或恒等变形进行严密证明。\\
  【解析计算】\\
  法一（特殊对称值法）：\\
  假设三角形为最对称的等边三角形，即 \(A = B = C = 60^\circ\)。\\
  此时：\\
  \(\tan A + \tan B + \tan C = \tan 60^\circ + \tan 60^\circ + \tan 60^\circ = \sqrt{3} + \sqrt{3} + \sqrt{3} = 3\sqrt{3}\)。\\
  因为这是一种合法可取到的状态，所以“最小值”必然不可能大于 \(3\sqrt{3}\)。\\
  据此可以排除 C (\(4\sqrt{3}\)) 和 D (\(6\sqrt{3}\))。\\
  由于锐角三角形中三个角可以无限逼近直角，此时正切值趋于正无穷，因此不可能是最大值，只能是最小值。在 A 和 B 之间，由于 \(3\sqrt{3} \approx 5.196\) 大于 \(2\sqrt{3} \approx 3.464\)，且等边三角形是极其特殊的临界状态，因此最小值锁定为 B。\\
  
  法二（正面严密证明）：\\
  因为 \(\triangle ABC\) 是锐角三角形，所以 \(A, B, C \in (0, \frac{\pi}{2})\)。\\
  在区间 \((0, \frac{\pi}{2})\) 上，函数 \(g(x) = \tan x\) 的二阶导数 \(g''(x) = 2\tan x \sec^2 x > 0\)，故 \(g(x)\) 为严格下凸函数。\\
  根据琴生不等式（Jensen's Inequality），有：\\
  \(\tan A + \tan B + \tan C \geqslant 3 \tan\left(\dfrac{A + B + C}{3}\right)\)。\\
  因为 \(A + B + C = \pi\)，所以：\\
  \(\tan A + \tan B + \tan C \geqslant 3 \tan\left(\dfrac{\pi}{3}\right) = 3\sqrt{3}\)。\\
  等号当且仅当 \(A = B = C = \frac{\pi}{3}\) 时成立。\\
  故最小值为 B。%
}

% -----------------------------------------------------
% Q3: 恒成立必要条件 (指对混合不等式)
% -----------------------------------------------------
\newcommand{\qThreeStem}{若对任意 \(x > 0\)，都有 \(\ln x \leqslant ax\)，则实数 \(a\) 的取值范围是\blank{1.5cm}}
\newcommand{\qThreeOptions}{\mcoptions[4]{\(a \geqslant 0\)}{\(a \geqslant 1/e\)}{\(a \geqslant 1\)}{\(a > 1/e\)}}
\newcommand{\qThreeAnswer}{B}
\newcommand{\qThreeSolution}{%
  【解题思路】\\
  面对“对任意 \(x\) 恒成立”求参数范围问题，先代入特殊的“舒服值”，往往能以极快的速度推出参数的“必要条件”，从而排除大部分选项，再通过导数证明其充分性。\\
  【解析计算】\\
  法一（特殊关键值法）：\\
  由于不等式中含有 \(\ln x\)，最舒服的代入值显然是使其为 1 或 0 的特殊值。\\
  令 \(x = e\)，代入不等式 \(\ln x \leqslant ax\) 得：\\
  \(\ln e \leqslant ae \implies 1 \leqslant ae \implies a \geqslant \dfrac{1}{e}\)。\\
  这直接给出了参数 \(a\) 必须满足的必要条件。结合选项，可以直接锁定 B（或者 D，但边界等号通常是可取的，故 B 概率极大）。\\
  
  法二（正面分离参数与导数法）：\\
  由于 \(x > 0\)，不等式 \(\ln x \leqslant ax\) 等价于：\\
  \(a \geqslant \dfrac{\ln x}{x}\)。\\
  设 \(g(x) = \dfrac{\ln x}{x}\)，求导数：\\
  \(g'(x) = \dfrac{\dfrac{1}{x} \cdot x - \ln x \cdot 1}{x^2} = \dfrac{1 - \ln x}{x^2}\)。\\
  令 \(g'(x) = 0 \implies x = e\)。\\
  当 \(0 < x < e\) 时，\(g'(x) > 0\)，\(g(x)\) 单调递增；\\
  当 \(x > e\) 时，\(g'(x) < 0\)，\(g(x)\) 单调递减。\\
  所以，当 \(x = e\) 时，\(g(x)\) 取得最大值，其最大值为 \(g(e) = \dfrac{1}{e}\)。\\
  为了使 \(a \geqslant g(x)\) 恒成立，必须满足 \(a \geqslant g(e) = \dfrac{1}{e}\)。\\
  故实数 \(a\) 的取值范围是 B。%
}

% -----------------------------------------------------
% Q4: 对称变换 (极值点正负对称)
% -----------------------------------------------------
\newcommand{\qFourStem}{设定义在 \(\mathbb{R}\) 上的函数 \(f(x)\) 满足 \(f(x) + f(-x) = 2\)，且当 \(x = 0\) 时，\(f(x)\) 取得极值。则下列结论中一定正确的是\blank{1.5cm}}
\newcommand{\qFourOptions}{\mcoptions[4]{\(x = 0\) 是 \(f(x)\) 的极大值点}{\(x = 0\) 是 \(f(x)\) 的极小值点}{\(f(0) > 1\)}{\(f'(0) = 0\)}}
\newcommand{\qFourAnswer}{D}
\newcommand{\qFourSolution}{%
  【解题思路】\\
  本题考查对称变换下的极值性质。\\
  第一步：分析题干条件中蕴含的对称性；\\
  第二步：利用正负函数对称变换（\(f(x) \to 2 - f(x)\)）对选项进行排它；\\
  第三步：利用可导函数极值的必要条件锁定正确答案。\\
  【解析计算】\\
  法一（对称变换排除法）：\\
  1. 题干条件 \(f(x) + f(-x) = 2\) 意味着函数 \(f(x)\) 的图象关于点 \((0, 1)\) 中心对称。因此在对称中心处必有 \(f(0) = 1\)，故排除 C。\\
  2. 若函数 \(f(x)\) 满足题干条件，则函数 \(h(x) = 2 - f(x)\) 也必然满足该条件（因为 \(h(x) + h(-x) = 2 - f(x) + 2 - f(-x) = 4 - (f(x)+f(-x)) = 2\)）。\\
  然而，若 \(x = 0\) 是 \(f(x)\) 的极大值点，那么在 \(h(x) = 2 - f(x)\) 中，\(x = 0\) 就会退化为极小值点。由于这两个函数在数学地位上完全等价，这就意味着我们不可能绝对断言 \(x=0\) 到底是“极大值点”还是“极小值点”。\\
  因此，A、B 选项可以通过对称变换直接予以排除。最后锁定 D。\\
  
  法二（正面分析）：\\
  对等式 \(f(x) + f(-x) = 2\) 两边关于 \(x\) 求导：\\
  \(f'(x) - f'(-x) = 0 \implies f'(x) = f'(-x)\)，即导函数 \(f'(x)\) 是偶函数。\\
  因为 \(x = 0\) 是可导函数 \(f(x)\) 的极值点，根据费马引理（Fermat's Theorem on stationary points），可导函数在极值点处的导数必须为 0。\\
  所以必有 \(f'(0) = 0\)。\\
  故一定正确的是 D。%
}

% -----------------------------------------------------
% Q5: 图像临界逼近 (指数函数切线题)
% -----------------------------------------------------
\newcommand{\qFiveStem}{若过点 \((a,b)\) 可以作曲线 \(y = e^x\) 的两条切线，则下列结论一定正确的是\blank{1.5cm} （多选）}
\newcommand{\qFiveOptions}{\mcoptions[4]{\(a > 0\)}{\(b > 0\)}{\(b < e^a\)}{\(b > a + 1\)}}
\newcommand{\qFiveAnswer}{BC}
\newcommand{\qFiveFigure}[1][1.0]{%
  \begin{tikzpicture}[scale=#1, >=stealth]
    % Draw grid
    \draw[color=gridcolor, thin] (-2.2,-0.5) grid (2.2,3.5);
    % Draw axes
    \draw[->, thick] (-2.3,0) -- (2.3,0) node[right] {\small \(x\)};
    \draw[->, thick] (0,-0.6) -- (0,3.6) node[above] {\small \(y\)};
    \draw (0,0) node[below left] {\small \(O\)};
    
    % Shaded region (0 < b < e^a)
    \fill[gray!15] plot[domain=-2.2:1.25] (\x, {exp(\x)}) -- (1.25, 0) -- (-2.2, 0) -- cycle;
    \node at (0.6, 0.4) [color=gray!70!black] {\small 2条切线区};
    
    % Draw curve y = e^x
    \draw[thick, black, domain=-2.2:1.25, samples=100] plot (\x, {exp(\x)});
    \node at (1.4, 3.2) {\small \(y=e^x\)};
    
    % Tangent line boundary y = x + 1 (at x=0)
    \draw[dashed, black!60] (-1.5, -0.5) -- (2.2, 3.2);
    \node[rotate=45] at (-0.6, 0.7) {\scriptsize \(y=x+1\)};
    
    % Draw point P(0.5, 1.0) and tangent lines
    \coordinate (P) at (0.5, 1.0);
    \fill[black] (P) circle (1.5pt) node[below right] {\scriptsize \(P(a,b)\)};
    
    % Tangents
    \draw[thin, blue, domain=-1.2:1.7] plot (\x, {exp(1.2)*(\x-1.2) + exp(1.2)});
    \draw[thin, blue, domain=-2.2:1.7] plot (\x, {exp(-0.83)*(\x+0.83) + exp(-0.83)});
  \end{tikzpicture}%
}
\newcommand{\qFiveSolution}{%
  【解题思路】\\
  利用图象特殊位置与切线临界性质进行判定。\\
  第一步：先画出指数函数 \(y = e^x\) 及其图象特征（下凸性）；\\
  第二步：代入几个特殊的点坐标，排除不正确选项；\\
  第三步：通过切线斜率方程建立严密方程解的个数关系。\\
  【解析计算】\\
  法一（特殊值排除法）：\\
  1. 取 \(a = -1\)，此时若要作两条切线，只要点在曲线下方且在 \(x\) 轴上方即可。比如取 \(P(-1, 0.2)\)。可以作两条切线。此时 \(a > 0\) 不成立，排除 A。\\
  2. 取 \(a = 0, b = 0.5\)，此时点在切线 \(y = x+1\)（在原点处的切线）下方，可以作两条切线。但此时 \(b > a+1 \implies 0.5 > 1\) 不成立，排除 D。\\
  通过简单特殊值排除，直接锁定正确答案为 B、C。\\
  
  法二（正面几何意义与方程解法）：\\
  设切点坐标为 \((t, e^t)\)。由于曲线在 \((t, e^t)\) 处的导数为 \(e^t\)，所以切线方程为：\\
  \(y - e^t = e^t(x - t) \implies y = e^t(x - t + 1)\)。\\
  因为切线过点 \((a, b)\)，代入切线方程得：\\
  \(b = e^t(a - t + 1)\)。\\
  过点 \((a,b)\) 可以作两条切线，等价于关于 \(t\) 的方程 \(b = e^t(a - t + 1)\) 有两个不同的实数根。\\
  设右端函数为 \(g(t) = e^t(a + 1 - t)\)，求导数：\\
  \(g'(t) = e^t(a + 1 - t) - e^t = e^t(a - t)\)。\\
  令 \(g'(t) = 0 \implies t = a\)。\\
  当 \(t < a\) 时，\(g'(t) > 0\)，\(g(t)\) 单调递增；\\
  当 \(t > a\) 时，\(g'(t) < 0\)，\(g(t)\) 单调递减。\\
  所以，当 \(t = a\) 时，\(g(t)\) 取得极大值，也是最大值，其最大值为 \(g(a) = e^a\)。\\
  同时，当 \(t \to -\infty\) 时，\(g(t) \to 0\)；当 \(t \to +\infty\) 时，\(g(t) \to -\infty\)。\\
  所以，要使得水平直线 \(y = b\) 与曲线 \(y = g(t)\) 有两个交点，必须满足：\\
  \(0 < b < g(a) \implies 0 < b < e^a\)。\\
  这直接给出了两个核心条件：\(b > 0\) 且 \(b < e^a\)。\\
  故一定正确的是 B、C。%
}

% -----------------------------------------------------
% Q6: 抽象函数一：余弦模型 (Q6A)
% -----------------------------------------------------
\newcommand{\qSixStem}{设函数 \(f(x)\) 满足 \(f(x+y) + f(x-y) = f(x)f(y)\)，且 \(f(0) = 2\)，\(f(1) = 1\)。求 \(f(2026)\) 的值。}
\newcommand{\qSixAnswer}{\(-1\)}
\newcommand{\qSixSolution}{%
  【解题思路】\\
  抽象函数题直接求解往往无从下手，其本质是“特殊模型法”：识别题目中隐藏的数学经典函数模型。\\
  第一步：观察结构特征：左边是 \(x+y\) 与 \(x-y\) 的和，右边是乘积，联想余弦的和差角公式；\\
  第二步：构造特殊余弦模型 \(f(x) = 2\cos kx\)，代入验证并确定参数 \(k\)；\\
  第三步：利用求出的模型计算目标值。\\
  【解析计算】\\
  法一（特殊模型法）：\\
  观察关系式：\(f(x+y) + f(x-y) = f(x)f(y)\)。\\
  我们熟知三角恒等式中：\\
  \(\cos(A+B) + \cos(A-B) = 2\cos A\cos B\)。\\
  由此，尝试构造特殊函数模型：\(f(x) = 2\cos kx\)。\\
  代入等式左边：\\
  \(f(x+y) + f(x-y) = 2\cos k(x+y) + 2\cos k(x-y) = 4\cos kx \cos ky\)。\\
  代入等式右边：\\
  \(f(x)f(y) = 2\cos kx \cdot 2\cos ky = 4\cos kx \cos ky\)。\\
  等式左右两端恒等，说明该模型非常成功！\\
  根据已知条件：\\
  \(f(0) = 2\cos 0 = 2\)（完全吻合）；\\
  \(f(1) = 2\cos k = 1 \implies \cos k = \dfrac{1}{2}\)。\\
  在第一象限内，可取最简单的 \(k = \dfrac{\pi}{3}\)。\\
  所以，我们得到具体的函数解析式模型：\\
  \(f(x) = 2\cos\left(\dfrac{\pi}{3}x\right)\)。\\
  由此计算：\\
  \(f(2026) = 2\cos\left(\dfrac{2026\pi}{3}\right)\)。\\
  因为 \(2026 = 6 \times 337 + 4\)，所以：\\
  \(f(2026) = 2\cos\left(674\pi + \dfrac{4\pi}{3}\right) = 2\cos\left(\dfrac{4\pi}{3}\right) = 2 \cdot \left(-\dfrac{1}{2}\right) = -1\)。\\
  
  法二（周期性分析法）：\\
  在原式中，令 \(y = 1\) 得：\\
  \(f(x+1) + f(x-1) = f(x)f(1) \implies f(x+1) + f(x-1) = f(x)\)。\\
  移项得：\(f(x+1) = f(x) - f(x-1)\)。\\
  由此可得：\\
  \(f(1) = 1\), \(f(0) = 2\)；\\
  \(f(2) = f(1) - f(0) = 1 - 2 = -1\)；\\
  \(f(3) = f(2) - f(1) = -1 - 1 = -2\)；\\
  \(f(4) = f(3) - f(2) = -2 - (-1) = -1\)；\\
  \(f(5) = f(4) - f(3) = -1 - (-2) = 1\)；\\
  \(f(6) = f(5) - f(4) = 1 - (-1) = 2\)；\\
  \(f(7) = f(6) - f(5) = 2 - 1 = 1\)。\\
  发现函数值呈周期性循环：\(2, 1, -1, -2, -1, 1, 2, 1, \dots\)，其周期为 \(T = 6\)。\\
  因为 \(2026 = 6 \times 337 + 4\)，所以：\\
  \(f(2026) = f(4) = -1\)。%
}

% -----------------------------------------------------
% Q7: 抽象函数二：对数模型 (Q6B)
% -----------------------------------------------------
\newcommand{\qSevenStem}{设对任意 \(x > 0\)，函数 \(f(x)\) 满足 \(f(xy) = f(x) + f(y) + 1\)，且 \(f(4) = 2\)。求 \(f\left(\dfrac{1}{2}\right)\) 的值。}
\newcommand{\qSevenAnswer}{\(-\dfrac{5}{2}\)}
\newcommand{\qSevenSolution}{%
  【解题思路】\\
  本题也是典型的抽象函数模型化。\\
  第一步：观察结构特征，左侧乘法拆成右侧加法，并且多了一个常数 1；\\
  第二步：通过加常数平移构造标准对数函数模型 \(g(x) = f(x) + 1\)；\\
  第三步：利用对数公式及其性质求解目标值。\\
  【解析计算】\\
  法一（特殊模型法）：\\
  已知关系式为：\(f(xy) = f(x) + f(y) + 1\)。\\
  为了消去常数 1，令 \(g(x) = f(x) + 1\)。代入等式：\\
  \(g(xy) - 1 = (g(x) - 1) + (g(y) - 1) + 1 \implies g(xy) = g(x) + g(y)\)。\\
  这是最经典的对数函数模型。尝试构造对数函数模型：\(g(x) = \log_k x\)。\\
  已知 \(f(4) = 2 \implies g(4) = f(4) + 1 = 3\)。\\
  所以 \(\log_k 4 = 3 \implies k^3 = 4 \implies k = 4^{1/3}\)。\\
  所以，具体的函数模型为：\\
  \(g(x) = \log_{4^{1/3}} x = 3\log_4 x\)。\\
  所以，具体的抽象函数解析式为：\\
  \(f(x) = g(x) - 1 = 3\log_4 x - 1\)。\\
  代入目标值 \(x = \dfrac{1}{2}\) 得：\\
  \(f\left(\dfrac{1}{2}\right) = 3\log_4\left(\dfrac{1}{2}\right) - 1 = 3 \cdot \left(-\dfrac{1}{2}\right) - 1 = -\dfrac{3}{2} - 1 = -\dfrac{5}{2}\)。\\
  
  法二（正面赋值计算法）：\\
  在原等式 \(f(xy) = f(x) + f(y) + 1\) 中：\\
  1. 令 \(x = y = 1\)，得 \(f(1) = f(1) + f(1) + 1 \implies f(1) = -1\)。\\
  2. 令 \(x = y = 2\)，得 \(f(4) = 2f(2) + 1\)。\\
  因为 \(f(4) = 2\)，所以 \(2 = 2f(2) + 1 \implies f(2) = \dfrac{1}{2}\)。\\
  3. 令 \(x = 2, y = \dfrac{1}{2}\)，得：\\
  \(f\left(2 \cdot \dfrac{1}{2}\right) = f(2) + f\left(\dfrac{1}{2}\right) + 1 \implies f(1) = f(2) + f\left(\dfrac{1}{2}\right) + 1\)。\\
  将已算出的 \(f(1) = -1\) 和 \(f(2) = \dfrac{1}{2}\) 代回：\\
  \(-1 = \dfrac{1}{2} + f\left(\dfrac{1}{2}\right) + 1 \implies -1 = \dfrac{3}{2} + f\left(\dfrac{1}{2}\right) \implies f\left(\dfrac{1}{2}\right) = -1 - \dfrac{3}{2} = -\dfrac{5}{2}\)。%
}

% -----------------------------------------------------
% Q8: 递推数列下界模型 (Q7)
% -----------------------------------------------------
\newcommand{\qEightStem}{已知数列 \(\{a_n\}\) 满足 \(a_1 = 1\)，\(a_2 = 1\)，且对任意正整数 \(n\)，都有 \(a_{n+2} > a_{n+1} + a_n\)，则下列一定成立的是\blank{1.5cm}}
\newcommand{\qEightOptions}{\mcoptions[4]{\(a_5 > 6\)}{\(a_5 > 5\)}{\(a_5 = 5\)}{\(a_5 < 5\)}}
\newcommand{\qEightAnswer}{B}
\newcommand{\qEightSolution}{%
  【解题思路】\\
  在处理含有不等号的递推数列时，我们常常需要寻找到它的“最低增长界限模型”，即把不等式退化为等式进行分析。\\
  第一步：将不等号退化为等号，构造最慢增长模型（即经典的斐波那契数列模型）；\\
  第二步：计算最慢模型在目标项的值；\\
  第三步：根据严格大于号，进行累加不等式推导，卡住范围并排除错误选项。\\
  【解析计算】\\
  法一（退化临界模型法）：\\
  我们将递推不等式 \(a_{n+2} > a_{n+1} + a_n\) 退化为最极端的临界等式：\\
  \(F_{n+2} = F_{n+1} + F_n\)，且 \(F_1 = 1, F_2 = 1\)。\\
  这是非常著名的斐波那契数列（Fibonacci Sequence）。\\
  我们计算前几项：\\
  \(F_1 = 1\),\\
  \(F_2 = 1\),\\
  \(F_3 = 1 + 1 = 2\),\\
  \(F_4 = 2 + 1 = 3\),\\
  \(F_5 = 3 + 2 = 5\)。\\
  因为实际递推关系为严格的“大于号”，故数列的增长速度必然严格大于斐波那契数列。\\
  由此可以合理猜测 \(a_5 > 5\)。锁定答案 B。\\
  同时因为只要增长多一点点（如 \(a_3 = 2.01\), \(a_4 = 3.02\), \(a_5 = 5.04\)）就符合题意，因此不可能保证 \(a_5 > 6\)，排除 A。\\
  
  法二（正面不等式推导）：\\
  由已知：\\
  \(a_1 = 1\)，\(a_2 = 1\)。\\
  因为 \(a_3 > a_2 + a_1\)，所以 \(a_3 > 1 + 1 = 2\)。\\
  因为 \(a_4 > a_3 + a_2\)，且 \(a_3 > 2, a_2 = 1\)，所以：\\
  \(a_4 > 2 + 1 = 3\)。\\
  因为 \(a_5 > a_4 + a_3\)，且 \(a_4 > 3, a_3 > 2\)，所以：\\
  \(a_5 > 3 + 2 = 5\)。\\
  故一定成立的是 B。%
}
